\documentclass[11pt,a4paper]{article}
\usepackage[margin=0.65in,top=0.58in,bottom=0.58in]{geometry}
\usepackage[T1]{fontenc}
\usepackage[utf8]{inputenc}
\usepackage{mathptmx,amsmath,booktabs,microtype,enumitem,xcolor,hyperref,fancyhdr,cite}
\hypersetup{colorlinks=true,linkcolor=black,citecolor=black,urlcolor=blue,
 pdftitle={Physical Consistency of Open Video Models: Project Milestone Report},
 pdfauthor={JIANG Huanyun}}
\definecolor{navy}{RGB}{25,48,73}
\pagestyle{fancy}\fancyhf{}
\fancyfoot[L]{\footnotesize AIST4010 \enspace | \enspace Project Milestone}
\fancyfoot[R]{\footnotesize\thepage\,/\,3}
\renewcommand{\headrulewidth}{0pt}
\setlength{\parindent}{0pt}\setlength{\parskip}{5pt}
\setlength{\tabcolsep}{7pt}
\newcommand{\sect}[2]{\vspace{4pt}{\large\bfseries\color{navy}#1. #2}\par\vspace{2pt}}
% Generated from archived experiment JSON; do not hand-edit.
\newcommand{\ParamCount}{29,465}
\newcommand{\TrainCount}{2,048}
\newcommand{\ValCount}{256}
\newcommand{\TestCount}{256}
\newcommand{\BestEpoch}{12}
\newcommand{\TestIoU}{0.970}
\newcommand{\TestError}{0.109}
\newcommand{\BaselineIoU}{0.901}
\newcommand{\BaselineError}{2.770}
\newcommand{\LossFirst}{1.448}
\newcommand{\ValIoUFirst}{0.000}
\newcommand{\LossThird}{0.529}
\newcommand{\ValIoUThird}{0.863}
\newcommand{\LossSixth}{0.043}
\newcommand{\ValIoUSixth}{0.947}
\newcommand{\LossLast}{0.022}
\newcommand{\ValIoULast}{0.974}
\newcommand{\WanGenerated}{12}
\newcommand{\WanTracked}{11}
\newcommand{\WanFitted}{0}
\newcommand{\CogGenerated}{12}
\newcommand{\CogTracked}{5}
\newcommand{\CogFitted}{0}
\newcommand{\FallError}{0.054}
\newcommand{\ProjectileError}{0.162}
\newcommand{\PeriodError}{0.138}

\begin{document}
\fontsize{10.5}{12.4}\selectfont
\begin{center}
{\LARGE\bfseries Physical Consistency of Open Video Models}\par\vspace{3pt}
{\large Project Milestone Report}\par\vspace{5pt}
\textbf{JIANG Huanyun (1155233165)}\quad AIST4010 Foundation of Applied Deep Learning
\end{center}
\vspace{-5pt}\hrule\vspace{5pt}

\sect{1}{Introduction}
Text-to-video generators can produce plausible images without reliably reproducing
the dynamics of the world. This distinction matters when generated videos are
used as synthetic experience for embodied agents: an attractive sequence may
still misrepresent falling objects, collisions, or support. Physical evaluation
therefore requires more than visual-quality scores. It also requires a trustworthy
measurement process: a poor tracker can make a correct trajectory look wrong,
while a visually convincing but unmeasurable scene cannot support a numerical
claim about gravity.

This project asks whether open-weight video models can generate quantitatively
consistent elementary motion under a reproducible, consumer-GPU evaluation
protocol. The overall plan has two parts. \textbf{Part A} extracts trajectories
from simple scenes and estimates acceleration, pendulum period, or deceleration.
\textbf{Part B}, still planned, evaluates visible violations in complex scenes
using a fixed, blinded vision-language-model (VLM) judging protocol. Synthetic
animations with known physics validate the numerical pipeline before it is used
on generated videos.

The expected contributions are a traceable local benchmark, explicit separation
of measurement failure from physical inconsistency, and a lightweight learned
detector that can be compared with classical segmentation. The original proposal
focused on evaluating frozen generators; this milestone adds supervised detector
training to address the observed measurement bottleneck. Current achievements
include both generator deployments, 24 audited clips, synthetic validation, and
an actually trained segmentation network. Reliable physical estimates on
generated clips and calibrated Part B results remain open goals.

\sect{2}{Related Work}
\textbf{Generative foundations.} Denoising diffusion probabilistic models
learn a reverse denoising process~\cite{Ho2020denoising}. Latent diffusion reduces
the computational burden by operating in an autoencoder's representation
space~\cite{Rombach2021high}, while diffusion transformers replace convolutional
denoisers with scalable latent-patch transformers~\cite{Peebles2022scalable}.
These advances help explain the feasibility of modern video generators, but
their synthesis objectives do not explicitly guarantee Newtonian dynamics.
Wan supplies an efficient 1.3B open-weight model~\cite{Wan2025open}; CogVideoX
uses a spatiotemporal VAE and expert transformer for text--video
fusion~\cite{Yang2024cogvideox}. They are the frozen generators evaluated here,
rather than models trained from scratch in this project.

\textbf{Quality versus physics.} VBench decomposes video quality into 16
dimensions, including temporal consistency and motion
smoothness~\cite{Huang2023vbench}. VideoPhy instead studies physical commonsense
in material interactions and evaluates \emph{both open and closed}
models~\cite{Bansal2024videophy}. PhyGenBench spans 160 prompts and 27 physical
laws, with a hierarchical VLM-based evaluator~\cite{Meng2024towards}. Physics-IQ
tests several physical domains and finds that realism and physical understanding
can dissociate~\cite{Motamed2025generative}. Together, these benchmarks motivate
physics-specific evaluation, while exposing the need to distinguish prompt
adherence, visible plausibility, and quantitatively measurable dynamics.

\textbf{Quantitative diagnosis and detection.} Kang et al. use controlled 2D
simulations to test physical-law generalization, finding case-based behavior and
out-of-distribution failures~\cite{Kang2024video}. Thus, the contribution here is
not the first quantitative physics benchmark; it is a small, auditable
implementation that reports the coverage and failure modes of its measurements.
For a trainable alternative to color rules, U-Net's encoder--decoder and skip
connections provide a localization-oriented segmentation
design~\cite{Ronneberger2015convolutional}. Adam supplies adaptive gradient-based
optimization~\cite{Kingma2014adam}. The auxiliary training pilot adapts these
established tools to synthetic ball masks; its purpose is better measurement,
not evidence that a generator has learned physics.

\newpage
\sect{3}{Data}
\textbf{Generated-video audit.} The completed dataset contains 24 clips:
two models, four English prompts, and three seeds (0, 1, 2). The phenomena are
free fall, horizontal projectile motion, a pendulum, and frictional horizontal
motion. Each prompt specifies a red ball of diameter 0.1 m; the pendulum string
is stated to be 1 m. Wan outputs 81 frames at 832$\times$480 and 16 fps;
CogVideoX outputs 49 frames at 720$\times$480 and 8 fps. Thus, the audit contains
1,560 frames in total. Videos, prompts, seeds, configuration, source hashes,
checkpoint hashes, and failure records are retained locally.

\textbf{Training and validation sources.} A separate procedural renderer supplies
2,560 RGB images at 96$\times$96 with exact binary ball masks:
\TrainCount{} training, \ValCount{} validation, and \TestCount{} test images,
using independent random-seed streams. Position, radius, shading, shadows,
thin red strings, blur, and noise vary. A typical example is a shaded red disk
with an attached string; only the disk belongs to the target mask. These are
synthetic images, not annotations of generator outputs. Four physics animations
provide independent kinematic checks. Forty-eight additional Part B videos and
four control-prompt types are planned; their evaluation is not yet complete.

\sect{4}{Approach}
\textbf{Frozen-model evaluation.} Both generators run through PyTorch/diffusers
with 50 sampling steps, guidance 5.0, and fixed seeds on separate RTX 5090 GPUs.
Classical red-component segmentation and temporal continuity recover ball
centroids. Nominal FPS supplies time; the prompt-stated diameter supplies a
conditional pixel-to-metre scale. Least-squares models include
\begin{equation*}
 y(t)=y_0+v_0t+\tfrac12gt^2,\qquad
 x(t)=x_0+A\sin(2\pi t/T+\phi),\qquad
 T_{\rm ref}=2\pi\sqrt{L/9.8}.
\end{equation*}
Projectile fits check both axes; horizontal fits also allow a stationary tail
after stopping. Acceptance requires $R^2\geq0.6$ and residual RMS $\leq3$ pixels,
reliable scale calibration, and sufficient temporal evidence. Missing targets,
internal disappearance, elongated components, or unstable apparent size are
recorded as failures. Rejected observations remain in the denominator.

\textbf{Trainable detector and baselines.} A U-Net-inspired network has
\ParamCount{} parameters, channel widths 8/16/32, two downsampling stages,
skip connections, and a one-channel mask head. It was trained for 12 epochs
with Adam, learning rate $10^{-3}$, batch size 64, and seed 7, minimizing binary
cross-entropy plus soft Dice loss. The checkpoint with highest validation mean
intersection-over-union (IoU) is selected before evaluating the held-out test
set. Masks use threshold 0.5. The immediate baseline is raw red-color
thresholding; comparisons with the full connected-component tracker and
background subtraction on labeled generated frames are planned.

\textbf{Evaluation and planned extension.} Detector metrics are mean mask IoU
and centroid error in pixels; absent predictions receive a full-image-diagonal
error penalty. Generator metrics are measurement coverage and, only when valid,
parameter error. Planned Part B uses eight sampled frames, opaque video IDs,
a fixed output schema, and manually verified calibration examples before
blinded judging. Static-frame judgments will not substitute for dynamical tests.

\vspace{8pt}
\begin{center}
\begin{tabular}{lll}
\toprule
\textbf{Component} & \textbf{Status at this milestone} & \textbf{Next validation target}\\
\midrule
Wan / CogVideoX & Deployed; weights frozen & Measurable physical trajectories\\
Classical detector & Synthetic gate passed & Robustness on generated scenes\\
Tiny U-Net & Trained; held-out synthetic test & Human-labeled generated frames\\
Part B judge & Protocol/prompt scaffolding & Calibration and blinded evaluation\\
\bottomrule
\end{tabular}
\end{center}

\newpage
\sect{5}{Preliminary Results}
\textbf{Actual deep-learning training.} Table~\ref{tab:training} records the
supervised segmentation run. Initial background-only predictions persist through
epoch 2; foreground masks emerge by epoch 3. Training loss and validation IoU
then improve together. The selected epoch is \BestEpoch{}. On 256 held-out
synthetic images, mean IoU is \TestIoU{} and centroid error is \TestError{} px,
versus \BaselineIoU{} and \BaselineError{} px for raw color thresholding.
This supports learning to separate the ball from the renderer's nuisance
features, but not transfer to generated video: no manually labeled real-output
test set has yet been evaluated. Results use one training seed.

\vspace{2pt}
\noindent\begin{minipage}[t]{0.47\textwidth}\centering
\small
\refstepcounter{table}\label{tab:training}
\textbf{Table \thetable. Training history}\par\vspace{3pt}
\begin{tabular}{rrr}
\toprule Epoch & Training loss & Validation IoU\\\midrule
1 & \LossFirst & \ValIoUFirst\\
3 & \LossThird & \ValIoUThird\\
6 & \LossSixth & \ValIoUSixth\\
12 & \LossLast & \ValIoULast\\\bottomrule
\end{tabular}
\end{minipage}\hfill
\begin{minipage}[t]{0.50\textwidth}\centering
\small
\refstepcounter{table}\label{tab:audit}
\textbf{Table \thetable. Frozen 24-clip audit}\par\vspace{3pt}
\begin{tabular}{lrrr}
\toprule Model & Generated & Tracked & Valid fit\\\midrule
Wan 1.3B & \WanGenerated & \WanTracked & \WanFitted\\
CogVideoX 5B & \CogGenerated & \CogTracked & \CogFitted\\\bottomrule
\end{tabular}\par\vspace{4pt}
{\footnotesize The trained detector has not re-scored this audit.}
\end{minipage}
\vspace{6pt}

\textbf{Physics audit and obstacles.} All 24 videos passed decoding and metadata
checks, but none passed the full measurement protocol (Table~\ref{tab:audit}):
eight failed tracking, nine failed calibration, and seven failed fitting.
Observed obstacles include truncated or extra objects, string-like components,
size changes, and non-fitting motion. Zero accepted fits means \emph{insufficient
reliable measurements}, not zero physical correctness. Synthetic relative errors
are approximately \FallError\% and \ProjectileError\% for free-fall and projectile
$g$, and \PeriodError\% for pendulum $T$; this does not eliminate domain shift.

\textbf{Next steps.} Annotate a held-out subset of generated frames; compare
learned and classical localization, then revalidate scale and trajectory fitting.
Preserve the current audit when revising prompts or detectors. Subsequently
complete Part B calibration and 48 additional clips. The immediate research
priority is valid measurement coverage, before a model-level physics ranking.

\vspace{1pt}
\begingroup
\fontsize{8.4}{9.6}\selectfont
\setlength{\parskip}{0pt}
\begin{thebibliography}{99}
\setlength{\itemsep}{1.8pt}\setlength{\parsep}{0pt}
% Generated from references.bib with verified venue enrichment.
\bibitem{Ho2020denoising} J. Ho et al., ``Denoising Diffusion Probabilistic Models,'' \emph{NeurIPS}, 2020.
\bibitem{Rombach2021high} R. Rombach et al., ``High-Resolution Image Synthesis with Latent Diffusion Models,'' \emph{CVPR, pp. 10684--10695}, 2022.
\bibitem{Peebles2022scalable} W. Peebles and S. Xie, ``Scalable Diffusion Models with Transformers,'' \emph{ICCV, pp. 4195--4205}, 2023.
\bibitem{Wan2025open} Team Wan et al., ``Wan: Open and Advanced Large-Scale Video Generative Models,'' \emph{arXiv:2503.20314}, 2025.
\bibitem{Yang2024cogvideox} Z. Yang et al., ``CogVideoX: Text-to-Video Diffusion Models with An Expert Transformer,'' \emph{ICLR}, 2025.
\bibitem{Huang2023vbench} Z. Huang et al., ``VBench: Comprehensive Benchmark Suite for Video Generative Models,'' \emph{CVPR, pp. 21807--21818}, 2024.
\bibitem{Bansal2024videophy} H. Bansal et al., ``VideoPhy: Evaluating Physical Commonsense for Video Generation,'' \emph{arXiv:2406.03520}, 2024.
\bibitem{Meng2024towards} F. Meng et al., ``Towards World Simulator: Crafting Physical Commonsense-Based Benchmark for Video Generation,'' \emph{arXiv:2410.05363}, 2024.
\bibitem{Motamed2025generative} S. Motamed et al., ``Do generative video models understand physical principles?'' \emph{WACV}, 2026.
\bibitem{Kang2024video} B. Kang et al., ``How Far is Video Generation from World Model: A Physical Law Perspective,'' \emph{ICML}, 2025.
\bibitem{Ronneberger2015convolutional} O. Ronneberger et al., ``U-Net: Convolutional Networks for Biomedical Image Segmentation,'' \emph{MICCAI, pp. 234--241}, 2015.
\bibitem{Kingma2014adam} D. P. Kingma and J. Ba, ``Adam: A Method for Stochastic Optimization,'' \emph{ICLR}, 2015.
\bibitem{Kassis2026scientific} T. Kassis et al., ``Scientific Agent Skills: A Library of Procedural Knowledge for Research Agents,'' \emph{arXiv:2609.00065}, 2026.

\end{thebibliography}
\endgroup
\vspace{-2pt}
{\fontsize{8}{9}\selectfont\textit{AI assistance:} OpenCode and Scientific Agent
Skills~\cite{Kassis2026scientific} assisted coding, drafting, and reference checks.
Training and evaluation were executed locally; planned work is identified explicitly.}
\end{document}
